\documentclass[aspectratio=169,11pt]{beamer}
\usepackage[utf8]{inputenc}
\usepackage[T1]{fontenc}
\usepackage[french]{babel}
\usepackage{helvet}
\renewcommand{\familydefault}{\sfdefault}
\usepackage{hyperref}
\definecolor{bleu}{HTML}{0000FF}
\definecolor{fond}{HTML}{F4F4EF}
\setbeamercolor{background canvas}{bg=fond}
\setbeamercolor{structure}{fg=bleu}
\setbeamercolor{title}{fg=bleu}
\setbeamercolor{frametitle}{fg=bleu,bg=fond}
\setbeamercolor{normal text}{fg=black}
\setbeamertemplate{navigation symbols}{}
\setbeamertemplate{itemize item}{\textcolor{bleu}{\tiny$\blacksquare$}}
\setbeamertemplate{footline}{\hfill\footnotesize\insertframenumber\hspace{1em}\vskip6pt}
\hypersetup{colorlinks=true,urlcolor=bleu,linkcolor=bleu}
\title{Fonds monétaire international}
\subtitle{Fonctionnement et données, vue d'ensemble}
\author{Veille institutionnelle}
\date{2026-11-03}

\begin{document}
\begin{frame}\titlepage\end{frame}

\section{Histoire et mission}
\begin{frame}{Histoire et mission (1/2)}
\begin{itemize}
\item Conçu à Bretton Woods en 1944, opérations en mars 1947
\item 1971 à 1973 : fin des parités fixes, flottement des monnaies
\item Crises de la dette (1980), Asie (1997), Argentine (2001), zone euro (2010)
\item 2008 : crise financière mondiale. 2020 : Covid-19
\end{itemize}
\end{frame}

\begin{frame}{Histoire et mission (2/2)}
\begin{itemize}
\item 2021 : allocation de DTS d'environ 650 milliards USD
\item Mission : stabilité monétaire et financière, coopération, croissance durable
\end{itemize}
\end{frame}

\begin{frame}{Histoire et mission : ressources à suivre}
\small
Veille hebdomadaire.
\begin{itemize}
\item \href{https://www.imf.org/en/About/Factsheets/Sheets/2023/IMF-at-a-glance}{Le FMI en bref}\\{\footnotesize Fiche institutionnelle, publication : annuelle}
\item \href{https://www.imf.org/en/About/Factsheets/Sheets/2016/08/01/14/45/The-IMF-and-the-World-Bank}{Histoire du FMI}\\{\footnotesize Fiche institutionnelle, publication : annuelle}
\item \href{https://www.imf.org/en/Countries}{Pays membres}\\{\footnotesize Page pays, publication : mensuelle}
\end{itemize}
\end{frame}

\section{Gouvernance}
\begin{frame}{Gouvernance (1/2)}
\begin{itemize}
\item Conseil des gouverneurs et CMFI de 24 membres, deux réunions par an
\item Conseil d'administration de 25 administrateurs
\item Quotes-parts : souscription, accès au financement, droits de vote
\item Droit de blocage des États-Unis sur les décisions à 85 \%
\end{itemize}
\end{frame}

\begin{frame}{Gouvernance (2/2)}
\begin{itemize}
\item 14e revue effective en 2016, 16e revue approuvée en décembre 2023
\item Directeur général par convention européen, premier adjoint américain
\end{itemize}
\end{frame}

\begin{frame}{Gouvernance : ressources à suivre}
\small
Veille hebdomadaire.
\begin{itemize}
\item \href{https://www.imf.org/en/About/Factsheets/Sheets/2016/07/14/12/21/Governance-Structure}{Gouvernance du FMI}\\{\footnotesize Fiche institutionnelle, publication : annuelle}
\item \href{https://www.imf.org/en/About/Factsheets/Sheets/2016/07/14/12/21/IMF-Quotas}{Quotes-parts}\\{\footnotesize Fiche institutionnelle, publication : annuelle}
\item \href{https://www.imf.org/en/About/Executive-Board}{Conseil d'administration}\\{\footnotesize Page institutionnelle, publication : mensuelle}
\item \href{https://www.imf.org/external/np/sec/memdir/members.aspx}{Membres et quotes-parts}\\{\footnotesize Liste, publication : annuelle}
\end{itemize}
\end{frame}

\section{Financement et ressources}
\begin{frame}{Financement et ressources (1/2)}
\begin{itemize}
\item Quotes-parts : ressource principale, environ 477 milliards de DTS
\item Accords d'emprunt NAB et accords bilatéraux
\item Capacité de prêt de l'ordre de 1 000 milliards USD
\item Revenus : intérêts, commissions, surcharges
\end{itemize}
\end{frame}

\begin{frame}{Financement et ressources (2/2)}
\begin{itemize}
\item Fonds fiduciaires : PRGT, CCRT, RST
\end{itemize}
\end{frame}

\begin{frame}{Financement et ressources : ressources à suivre}
\small
Veille hebdomadaire.
\begin{itemize}
\item \href{https://www.imf.org/en/About/Factsheets/Sheets/2023/Where-the-IMF-gets-its-money}{Origine des ressources}\\{\footnotesize Fiche institutionnelle, publication : annuelle}
\item \href{https://www.imf.org/external/np/fin/tad/extarr2.aspx}{Situation financière}\\{\footnotesize Données financières, publication : mensuelle}
\end{itemize}
\end{frame}

\section{Surveillance}
\begin{frame}{Surveillance (1/2)}
\begin{itemize}
\item Base juridique : article IV des statuts
\item Consultations article IV en général annuelles, rapport publié avec accord du pays
\item Surveillance multilatérale : WEO, GFSR, Fiscal Monitor, External Sector Report
\item FSAP depuis 1999 : évaluation de la stabilité financière
\end{itemize}
\end{frame}

\begin{frame}{Surveillance (2/2)}
\begin{itemize}
\item Revue complète de la surveillance tous les cinq ans environ
\end{itemize}
\end{frame}

\begin{frame}{Surveillance : ressources à suivre}
\small
Veille hebdomadaire.
\begin{itemize}
\item \href{https://www.imf.org/en/About/Factsheets/Sheets/2023/IMF-Surveillance}{Surveillance du FMI}\\{\footnotesize Fiche institutionnelle, publication : annuelle}
\item \href{https://www.imf.org/en/Publications/CR}{Rapports article IV}\\{\footnotesize Base de rapports, publication : continue}
\item \href{https://www.imf.org/en/Publications/FSAP}{FSAP}\\{\footnotesize Base de rapports, publication : trimestrielle}
\end{itemize}
\end{frame}

\section{Prêts et programmes}
\begin{frame}{Prêts et programmes (1/2)}
\begin{itemize}
\item SBA et EFF pour les besoins de balance des paiements
\item RFI et RCF : financements d'urgence
\item FCL, PLL, SLL : lignes de précaution
\item PRGT : prêts concessionnels. RSF : résilience climatique
\end{itemize}
\end{frame}

\begin{frame}{Prêts et programmes (2/2)}
\begin{itemize}
\item Conditionnalité : actions préalables, critères de réalisation, repères structurels
\item Dette : cadre commun du G20, table ronde mondiale, analyse de viabilité
\end{itemize}
\end{frame}

\begin{frame}{Prêts et programmes : ressources à suivre}
\small
Veille hebdomadaire.
\begin{itemize}
\item \href{https://www.imf.org/en/About/Factsheets/Sheets/2023/IMF-Lending}{Prêts du FMI}\\{\footnotesize Fiche institutionnelle, publication : annuelle}
\item \href{https://www.imf.org/en/Countries/ResRep}{Programmes en cours}\\{\footnotesize Page pays, publication : mensuelle}
\item \href{https://www.imf.org/en/Topics/imf-and-debt}{Dette et FMI}\\{\footnotesize Page thématique, publication : trimestrielle}
\end{itemize}
\end{frame}

\section{Droits de tirage spéciaux}
\begin{frame}{Droits de tirage spéciaux (1/2)}
\begin{itemize}
\item Créés en 1969 comme actif de réserve complémentaire
\item Panier : dollar, euro, yuan, yen, livre sterling
\item Allocation d'août 2021 : environ 456 milliards de DTS
\item Échange contre devises, recyclage vers les pays pauvres
\end{itemize}
\end{frame}

\begin{frame}{Droits de tirage spéciaux (2/2)}
\begin{itemize}
\item Révision du panier tous les cinq ans
\end{itemize}
\end{frame}

\begin{frame}{Droits de tirage spéciaux : ressources à suivre}
\small
Veille hebdomadaire.
\begin{itemize}
\item \href{https://www.imf.org/en/About/Factsheets/Sheets/2023/special-drawing-rights-sdr}{DTS}\\{\footnotesize Fiche institutionnelle, publication : annuelle}
\item \href{https://www.imf.org/external/np/fin/data/rms_sdrv.aspx}{Taux de DTS}\\{\footnotesize Données, publication : mensuelle}
\end{itemize}
\end{frame}

\section{Développement des capacités}
\begin{frame}{Développement des capacités}
\begin{itemize}
\item Assistance technique et formation, souvent gratuites pour les pays pauvres
\item Domaines : finances publiques, monnaie, supervision, statistiques
\item Centres régionaux : AFRITAC, CARTAC, Pacifique, Moyen-Orient
\item Cours en ligne et en présentiel
\end{itemize}
\end{frame}

\begin{frame}{Développement des capacités : ressources à suivre}
\small
Veille hebdomadaire.
\begin{itemize}
\item \href{https://www.imf.org/en/Capacity-Development}{Développement des capacités}\\{\footnotesize Page thématique, publication : trimestrielle}
\item \href{https://www.imf.org/en/Capacity-Development/training}{Formation}\\{\footnotesize Catalogue de cours, publication : trimestrielle}
\end{itemize}
\end{frame}

\section{Recherche et publications}
\begin{frame}{Recherche et publications (1/2)}
\begin{itemize}
\item WEO : avril et octobre, mises à jour en janvier et juillet
\item GFSR et Fiscal Monitor : avril et octobre
\item External Sector Report : annuel
\item Perspectives économiques régionales
\end{itemize}
\end{frame}

\begin{frame}{Recherche et publications (2/2)}
\begin{itemize}
\item Working Papers en flux continu
\end{itemize}
\end{frame}

\begin{frame}{Recherche et publications : ressources à suivre}
\small
Veille hebdomadaire.
\begin{itemize}
\item \href{https://www.imf.org/en/Publications/WEO}{World Economic Outlook}\\{\footnotesize Publication, publication : semestrielle}
\item \href{https://www.imf.org/en/Publications/GFSR}{Global Financial Stability Report}\\{\footnotesize Publication, publication : semestrielle}
\item \href{https://www.imf.org/en/Publications/FM}{Fiscal Monitor}\\{\footnotesize Publication, publication : semestrielle}
\item \href{https://www.imf.org/en/Publications/WP}{Working Papers}\\{\footnotesize Publication, publication : continue}
\end{itemize}
\end{frame}

\section{Données}
\begin{frame}{Données (1/2)}
\begin{itemize}
\item data.imf.org : portail principal, API SDMX sur api.imf.org
\item DataMapper : API JSON simple, historique et prévisions WEO
\item Bases : IFS, BOP, DOTS, GFS, COFER, CPI, FSI, PCPS
\item Normes de diffusion : e-GDDS, SDDS, SDDS Plus
\end{itemize}
\end{frame}

\begin{frame}{Données (2/2)}
\begin{itemize}
\item Les pays déclarent, le département des statistiques contrôle et publie
\item Ancienne API dataservices.imf.org hors service
\end{itemize}
\end{frame}

\begin{frame}{Données : ressources à suivre}
\small
Veille hebdomadaire.
\begin{itemize}
\item \href{https://data.imf.org}{Portail des données}\\{\footnotesize Portail, publication : mensuelle}
\item \href{https://www.imf.org/external/datamapper}{DataMapper}\\{\footnotesize Portail, publication : mensuelle}
\item \href{https://api.imf.org/external/sdmx/2.1/dataflow}{API SDMX}\\{\footnotesize API, publication : trimestrielle}
\item \href{https://dsbb.imf.org}{Tableau d'affichage DSBB}\\{\footnotesize Normes de diffusion, publication : trimestrielle}
\item \href{https://www.imf.org/en/Data}{Données du FMI}\\{\footnotesize Page thématique, publication : mensuelle}
\end{itemize}
\end{frame}

\subsection{data.imf.org : portail}
\begin{frame}{data.imf.org : portail}
\begin{itemize}
\item Portail principal des données statistiques du FMI
\item Explorateur de données, jeux de données, calendrier de diffusion
\item S'appuie sur l'API SDMX de api.imf.org
\end{itemize}
\end{frame}

\begin{frame}{data.imf.org : portail : ressources à suivre}
\small
Veille hebdomadaire.
\begin{itemize}
\item \href{https://data.imf.org}{Accueil du portail}\\{\footnotesize Portail, publication : hebdomadaire}
\item \href{https://data.imf.org/en/Datasets}{Jeux de données}\\{\footnotesize Médiathèque, publication : mensuelle}
\item \href{https://data.imf.org/en/Data-Explorer}{Explorateur de données}\\{\footnotesize Requêteur, publication : mensuelle}
\item \href{https://data.imf.org/en/news}{Actualités du portail}\\{\footnotesize Médiathèque, publication : hebdomadaire}
\end{itemize}
\end{frame}

\subsection{DataMapper}
\begin{frame}{DataMapper}
\begin{itemize}
\item Cartes et séries WEO avec une API JSON simple
\item Mélange historique et prévisions : lire l'année de coupure
\item Idéal pour des indicateurs macroéconomiques standards
\end{itemize}
\end{frame}

\begin{frame}{DataMapper : ressources à suivre}
\small
Veille hebdomadaire.
\begin{itemize}
\item \href{https://www.imf.org/external/datamapper}{DataMapper}\\{\footnotesize Portail, publication : mensuelle}
\item \href{https://www.imf.org/external/datamapper/api/v1/indicators}{Indicateurs DataMapper (API)}\\{\footnotesize API, publication : trimestrielle}
\end{itemize}
\end{frame}

\subsection{API SDMX}
\begin{frame}{API SDMX}
\begin{itemize}
\item Catalogue des dataflows : identifiant, version, agence
\item XML par défaut, JSON avec l'en-tête Accept
\item Ancienne API dataservices.imf.org retirée
\end{itemize}
\end{frame}

\begin{frame}{API SDMX : ressources à suivre}
\small
Veille hebdomadaire.
\begin{itemize}
\item \href{https://api.imf.org/external/sdmx/2.1/dataflow}{Liste des dataflows}\\{\footnotesize API, publication : hebdomadaire}
\end{itemize}
\end{frame}

\subsection{Normes de diffusion (DSBB)}
\begin{frame}{Normes de diffusion (DSBB)}
\begin{itemize}
\item e-GDDS, SDDS et SDDS Plus : engagements de publication des pays
\item Calendriers préannoncés et métadonnées
\item Utile pour savoir quand un pays publie
\end{itemize}
\end{frame}

\begin{frame}{Normes de diffusion (DSBB) : ressources à suivre}
\small
Veille hebdomadaire.
\begin{itemize}
\item \href{https://dsbb.imf.org}{Tableau d'affichage DSBB}\\{\footnotesize Normes de diffusion, publication : mensuelle}
\end{itemize}
\end{frame}

\subsection{Médiathèque : blogs, vidéos, podcasts}
\begin{frame}{Médiathèque : blogs, vidéos, podcasts}
\begin{itemize}
\item Blogs et notes de politique économique
\item Vidéos et podcasts de la direction et des économistes
\item Page des données du site institutionnel
\item À consulter chaque semaine pour les annonces
\end{itemize}
\end{frame}

\begin{frame}{Médiathèque : blogs, vidéos, podcasts : ressources à suivre}
\small
Veille hebdomadaire.
\begin{itemize}
\item \href{https://www.imf.org/en/Blogs}{Blogs du FMI}\\{\footnotesize Blog, publication : hebdomadaire}
\item \href{https://www.imf.org/en/Videos}{Vidéos}\\{\footnotesize Vidéo, publication : hebdomadaire}
\item \href{https://www.imf.org/en/Podcasts}{Podcasts}\\{\footnotesize Podcast, publication : hebdomadaire}
\item \href{https://www.imf.org/en/Data}{Page données du FMI}\\{\footnotesize Page thématique, publication : mensuelle}
\item \href{https://www.imf.org/en/Publications/WEO/weo-database}{Base de données WEO}\\{\footnotesize Base de données, publication : semestrielle}
\end{itemize}
\end{frame}

\section{Évaluation et transparence}
\begin{frame}{Évaluation et transparence}
\begin{itemize}
\item Bureau indépendant d'évaluation créé en 2001
\item Politique de transparence : publication des rapports avec l'accord des pays
\item Revues des programmes, des instruments, de la conditionnalité
\item Accès aux archives historiques après un délai
\end{itemize}
\end{frame}

\begin{frame}{Évaluation et transparence : ressources à suivre}
\small
Veille hebdomadaire.
\begin{itemize}
\item \href{https://ieo.imf.org}{Bureau indépendant d'évaluation}\\{\footnotesize Site d'évaluation, publication : trimestrielle}
\item \href{https://www.imf.org/en/About/Archives}{Archives du FMI}\\{\footnotesize Archives, publication : ponctuelle}
\end{itemize}
\end{frame}

\section{Plans stratégiques}
\begin{frame}{Plans stratégiques (1/2)}
\begin{itemize}
\item Programme d'orientation mondiale à chaque réunion de printemps et annuelle
\item Stratégie climat de 2021 et facilité pour la résilience
\item Stratégies de développement des capacités, genre, données
\item Revues de la surveillance et de la conditionnalité
\end{itemize}
\end{frame}

\begin{frame}{Plans stratégiques (2/2)}
\begin{itemize}
\item Dette souveraine : cadre commun du G20, table ronde mondiale
\item À vérifier sur imf.org avant citation
\end{itemize}
\end{frame}

\begin{frame}{Plans stratégiques : ressources à suivre}
\small
Veille hebdomadaire.
\begin{itemize}
\item \href{https://www.imf.org/en/Publications/Policy-Papers}{Documents de politique}\\{\footnotesize Base documentaire, publication : mensuelle}
\item \href{https://www.imf.org/en/Topics/climate-change}{Climat et FMI}\\{\footnotesize Page thématique, publication : trimestrielle}
\item \href{https://www.imf.org/en/Meetings}{Réunions de printemps et annuelles}\\{\footnotesize Calendrier, publication : semestrielle}
\end{itemize}
\end{frame}

\section{Critiques et débats}
\begin{frame}{Critiques et débats (1/2)}
\begin{itemize}
\item Conditionnalité et austérité des années 1980 et 1990
\item Crise asiatique de 1997 : recommandations jugées trop restrictives
\item Argentine et Grèce : erreurs reconnues, multiplicateurs budgétaires
\item Surcharges : coût supplémentaire pour les grands emprunteurs, réforme en 2024
\end{itemize}
\end{frame}

\begin{frame}{Critiques et débats (2/2)}
\begin{itemize}
\item Représentation : poids des pays avancés dans les quotes-parts
\end{itemize}
\end{frame}

\begin{frame}{Critiques et débats : ressources à suivre}
\small
Veille hebdomadaire.
\begin{itemize}
\item \href{https://ieo.imf.org}{Évaluations indépendantes}\\{\footnotesize Site d'évaluation, publication : trimestrielle}
\item \href{https://www.imf.org/en/Publications/Policy-Papers}{Documents de politique}\\{\footnotesize Base documentaire, publication : mensuelle}
\end{itemize}
\end{frame}

\end{document}
